\clearpage
\thispagestyle{empty}
\begin{abstract}
I denne specialeafhandling ser vi på problemet med at finde en approksimation af vektoren \(x\) fra det lineære ligningssystem \(Ax=b\), med hovedformålet at forkorte køretiden af den iterative metode conjugate gradient. Det gøres ved at transformere systemet med en forbetingelsesmatrix. Vores mål er at finde ud af, om det er muligt at lære denne matrix således, at den kan reducere køretiden for et generelt sæt af ukendte lineære systemer. Vores interesse for dette kommer fra studiet af kvantekromodynamik (QCD) med modellen Gitter-QCD, hvor den beregningsmæssige dominerende del af modellen er applikationen af den inverse Dirac-operatør, \(D\). Konvergeringsraten af conjugate gradient for fejlen i approksimationen af \(x\) kan blive afgrænset af betingelsesnummeret af matrixen, \(\kappa(A)\). Derfor: hvis det lineære system bliver transformeret af forbetingelsesmatrixen således at betingelsesnummeret reduceres, så forkortes køretiden af conjugate gradient.
\\

Undersøgelsen af læringen af forbetingelsesmatrixen bestod af to kategorier: strukturen af matrixen og læringsalgoritmen. Strukturen af forbetingelsesmatrixen startede som en øvre bidiagonalmatrix, med en identitet hoveddiagonal og en partitioneret superdiagonal med en koefficient for hver partitionering. Læringsalgoritmen virker ved at ændre en af koefficienterne i superdiagonalen, hvorefter den evaluerer med conjugate gradient for at bestemme om ændringen reducerer antallet af iterationer yderligere.
\\

For at bestemme hvornår en ændring af forbetingelsesmatrixen er en forbedring, testede vi læringsalgoritmen med tre forskellige forbedringsfunktioner: MEAN-forbedringsfunktionen, som sammenligner gennemsnittet af iterationstallene, MEDIAN-forbedringsfunktionen, som sammenligner medianen af ite-rationstallene og SIGN-forbedringsfunktionen, som sammenligner de individuelle lineære systemers ite-rationstal. Resultaterne fra evalueringen af de lærte forbetingelsesmatrixer viste, at SIGN-forbedrings-funktionen frembragte den bedste forbetingelsesmatrix i læringen. Men i en sammenligning med en kendt forbetingelsesmatrix, Jacobi forbetingelse, der bruger hoveddiagonalen af matrixen fra det lineære sy-stem som forbetingelsesmatrix, var iterationstallene af de forbetingede lineære systemer noget større og havde dermed en længere køretid. Derfor ændrede vi hoveddiagonalen af vores forbetingelsesmatrix fra identitetsmatrixen til den inverse af hoveddiagonalen af det lineære systems matrix. Her blev læringen og evalueringen noget bedre end vores tidligere forbetingelsesmatrixer, men dog kun med en marginal forskel til Jacobi forbetingelse.
\\

De næste strukturelle ændringer, vi testede, involverede superdiagonalen: mængden af partitione-ringer og antallet af ikke-nul superdiagonaler. Resultaterne af disse læringstests viste at ændringerne af superdiagonalerne ikke medførte en forbedring af forbetingelsesmatrixen. Men vi kunne dog konkludere, at kun at bruge den første superdiagonal og kun én koefficient til hele superdiagonalen gav den hurtigste læring af forbetingelsesmatrixen, uden at den blev forringet.
\\

Den sidste store ændring af forbetingelsesmatrixen var at tilføje en ikke-nul subdiagonal. Dette betyder at strukturen af forbetingelsesmatrixen nu er en tridiagonalmatrix med en hoveddiagonal som Jacobi forbetingelse og en partitioneret sub- og superdiagonal. Evalueringen af de lærte forbetingelsesmatrixer med den nye struktur viste en reduktion af gennemsnits iterationstallet på to gange så meget, som Jacobi forbetingelse reducerer. Efter undersøgelsen af den nye forbetingelsesmatrix, er vi kommet frem til, at den tridiagonale struktur med én koefficient i både sub- og superdiagonalerne giver de bedste resultater i evalueringen af de lærte forbetingelsesmatrixer. Vi kan hermed konkludere, at det er muligt at lære en generel forbetingelsesmatrix, der kan reducere antallet af iterationer, som en iterativ metode har brug for, for at finde en approksimation til et ukendt lineært ligningssystem. 
\end{abstract}
\clearpage


% I denne specialafhandling ser vi problemet med at finde approksimation af vektoren \(x\) i det standard lineære system \(Ax=b\), med hoved formålet at forkorte køretiden af den iterative metode konjugeret gradient. Det kan gøres ved at transformere system med en forbetingelsesmatrix. Vores mål at finde ud af om det er muligt at lære denne matrice sådan at den kan reducere køretiden for er genrealt sæt af ligende lineære systemer. Vores *** kommer fra studiet af Kvantekromodynamik (QCD) med modelen Gitter-QCD, hvor den beregningsmæssige dominerne del af modelen er applikationen af den inverse Dirac-operatør, \(D\). Konveregning raten af konjugeret gradient for fejlen i approksimationen af \(x\) kan blive afgrænset af betingelsesnummer af matricen, \(\kappa(A)\). Derfor hvis det lineære system bliver transformeret af forbetingelsesmatrixen sådan at betingelsesnummeret reduceres, forkortes køretiden af konjugeret gradient.
% \[\]
% I udforskiningen af lærering af en forbetingelsesmatrix betod af to kategorier, strukturen af matrixen og lærerings algoritmen. For strukturen af forbetingelsesmatrixen den startede som en øvre bidiagonalmatrix, med en identitet hoveddiagonal og en partitioenret superdiagonal med en koefficient for vær partitonering. Lærering algoritmen fungere ved at ændre en af koefficienterne i superdiagonalen løser de forbetingede lineær systemer for at for et iterationstal til at sammenligne med iterationstallet fra den tidligere forbetingede lineær systemer. 
% \[\]
% For at finde hvilke mål af iterationstallene producer den bedste lærte forbetingelsesmatrix testede vi lærering algoritmen med tre forskellige forbedringsfunktion: MEAN forbedringsfunktion sammenligner gennemsnits iterationstallet, MEDIAN forbedringsfunktion sammenligner median iterationstallet og SIGN forbedringsfunktion som sammenligner de individuale lineære systemers iterationstal. Resultaterne fra evaluering af de lærete forbetingelsesmatrixer var SIGN forbedringsfunktion bedste til at lære en forbetingelsesmatrix. I en sammenligning med et kendt forbetingelsesmatrix, Jacobi forbetingelse som bruger hoveddiagonalen af det lineær systems matrix som forbetingelsesmatrix, var resultaterne markant dårligere. Derfor prøvede vi at ændre hoveddiagonalen af vores forbetingelsesmatrix fra identiteten til den inverse af hoveddiagonalen af det lineær systems matrix. Her var læreringen og evaluering markant bedre end de tidligere forbetingelsesmatrixer, men ikke med nogen stor forskel til Jacobi forbetingelse.
% \[\]
% De næste strukturale til ændreniger involverede superdiagonalen, mængden af partioenringer og antalet af ikke-nul superdiagonaler. Fra de lærerings tests blev der ikke lærte en forbetingelsesmatrix som præterete bedre, men konkluderede at kun at bruge den første superdiagonal og kun en koefficient til hele superdiagonalen gave den hurtiges lærering af forbetingelsesmatrixen. Den sidste store ændring af forbetingelsesmatrixen var at have en ikke-nul subdiagonal, det betyder at strukturen af forbetingelsesmatrixen er en tridiagonalmatrix med en hoveddiagonal som Jacobi forbetingelse og en partitioenret sub- og superdiagonal. Ved evaluering af denne nye matrix reducerede den i gennemsnit iterationstallet af konjugeret gradient cirka to gange så meget som Jacobi forbetingelse. Derefter en mindre udforskiningen af den nye forbetingelsesmatrix er vi kommet frem til at der er muligt at lære en general forbetingelsesmatrix som præterere godt på ukendte lineær systemer.
% Conjugate gradient
% eng
% In this master thesis we will look at problem of inversion of matrices in the context of finding an approximation for the value of \(x\) in the linear system \(Ax=b\), with the main goal of decreasing the computational time of getting the approximation. This is be done with a preconditioning matrix, which is an approximation of the inverse of the original matrix. Our goal is to see if it is possible to learn this  preconditioning matrix such that it can decrease the computational cost of a set of linear systems. Our interest comes from the study of Quantum Chromodynamics (QCD) with the model Lattice QCD, where the computational dominate part of the model is the application of the inverse of the Dirac operator, \(D\).
% One such method to finding an approximation that we will derive is the Conjugate Gradient. The convergence rate for the error of the approximation  of this method can be found to conditioned on the condition number of the systems matrix \(\kappa(A)\). Therefore if the system is transformed such that the condition number is smaller, Conjugate Gradient will have a faster rate of convergence. This is done by the preconditioning matrix.
% The exploration of learning a preconditioner consisted of two main categories, the structure of the preconditioner and the learning algorithm. The structure of our preconditioner started as upper bidiagonal matrix with identity diagonal and a partitioned super diagonal. The learning algorithm would change a coefficient of from the super diagonal and compare the iteration count of Conjugate Gradient with the previous preconditioned system. By a test of four standard test matrices we found that with this preconditioner and learning algorithm, it can improve the preconditioner from an initial state to a state with fewer iteration count for the solver. The next major test involved the so called improvement functions, which determines when a change made to a coefficient improved the preconditioner. This was needed because when learning a preconditioner for a large set of linear systems you need one point of comparison between preconditioners. The three improvement functions we tested was the MEAN improvement function, which compared the mean solver iteration count; the MEDIAN improvement function, which compared the median solver iteration count; the SIGN improvement function, which does an comparison between the individual linear systems and count if a majorty of the systems have a lower iteration count.
% The results from this test was the SIGN improvement function learned the best preconditioner of the three but in a comparison with a simple known preconditioner, Jacobi preconditioning, the learned preconditioners preformed poorly compared to it. Therefore we changed diagonal of the preconditioner to the inverse of the linear system, which is the Jacobi preconditioner. The learned preconditioners with this structure we got strict improvement to the previous and results on par with Jacobi preconditioning.
% The next structural changes involved the superdiagonal, those being the number of partition and the number of nonzero superdiagonal. From this test we concluded that one nonzero superdiagonal and one partitioned resolved in the fastest learning and best preconditioners, but still on par with Jacobi. The final major change was adding nonzero values to the subdiagonal. By this change the learned preconditioners from the test preformed twice as well compared to Jacobi preconditioning. By the whole exploration the final conclusion of the this thesis is that learning a general preconditioner to work on a general set of similar matrices does work.
% unused
% The learning algorithm is a simple procedure that will accept the change made to the coefficients of the partitioned super diagonal based on the Conjugate Gradient iteration count of the preconditioned linear systems. 
% The exploration of the learning algorithm we tested three different parts. The first was about the changes made to the coefficients of the preconditioner. Where we compared a larger amount of changes against fewer but a more thorough exploration of the change, such that the total number of possible changes is equal. From this test 
% The first major test was to determine the measurement of improvement, a function to determine if a change made to a coefficient of the preconditioner improved it. Here we tested three different measures, the mean iteration count, the median iteration count, and the number linear systems that have a lower iteration count compared with the previous learning step. 
% Therefore we changed the diagonal of our preconditioner to the inverse of the diagonal of the systems matrix, similar to Jacobi preconditioning. This gave us learned preconditioners that preformed on par with Jacobi preconditioning. After we tested changes to the preconditioner: the number of partition, the number of nonzero upper diagonals, and lastly using both the sub and super diagonal. Where the two former 


% The results we got from the first test with this preconditioner is that it is possible to learn the coefficients of the upper diagonal for one matrix. The next part was learning a preconditioner that can work for a set of matrices. The data set we choose is a set of 2-dimensional discrete laplace operator, \(\), with perturbations on each nonzero entry to make every matrix in the set unique


% This method constructs the approximation for \(x\) from linear combination of an orthogonal basis from an increasingly larger subspace. It therefore in exact arithmetic uses a maximum number of iterations equal to the dimension of the system.

% The problem is the construction of a good preconditioning matrix is difficult. Because the condition number has to be reduced enough such that the extra computational cost of using a preconditioner does not undo the the decrease in iteration count for Conjugate Gradient. This is increased by the need to precondition a set of linear systems. 

% Second it is often unique to the matrix. Therefore if the need to precondition a larger set of linear systems a unique preconditioner is needed for each system. 
